\documentclass{article}
\usepackage[UTF8]{ctex}
\usepackage{amsmath,amssymb,amsfonts}
\usepackage{enumitem}
\usepackage{geometry}
\geometry{a4paper, margin=1in}
% === 错题本排版优化（不参与题目标识） ===
\setlength{\parindent}{0pt}
\setlength{\parskip}{0.35em plus 0.10em minus 0.08em}
\setlength{\emergencystretch}{3em}
\setlist{itemsep=0.25em, topsep=0.40em, parsep=0pt}
\allowdisplaybreaks[2]
% === 错题本排版优化结束 ===
\begin{document}
\section*{中山大学2025年数学分析真题}
\section*{153．中山大学2025年数学分析试卷}

\subsection{题目1}
1．填空题．每题5分，共50分．

\begin{itemize}
  \item[（1）] \(\lim _{n \rightarrow \infty} \sum_{k=1}^{n} \frac{\sin \frac{k}{n} \pi}{n+\frac{k}{n}}=\)
  \item[（2）] 设 \(f(x)\) 在 \(x=0\) 处可导，且 \(f(0)=2\) ，则 \(\lim _{n \rightarrow \infty}\left(\frac{f\left(\frac{2}{n}\right)}{f(0)}\right)^{n}=\) \(\_\_\_\_\)．
  \item[（3）] \(\sum_{n=1}^{\infty} \frac{n(n+1)}{2^{n}}=\)
  \item[（4）] \(\sum_{n=1}^{\infty} \frac{1}{3^{n} \sqrt{n+1} x^{3 n}}\) 的收敛域为 \(\_\_\_\_\)．
  \item[（5）] \(f(s, t)=e^{2 s t} \sin \left(s^{2}+t\right)\) ，则 \(\frac{\partial f}{\partial s}=\) \(\_\_\_\_\)．
  \item[（6）] \(f(x, y, z)=x^{2}+y^{2}+z^{2}\) ，则 \(\operatorname{div}(\nabla f)=\) \(\_\_\_\_\)．
  \item[（7）] 设 \(f(x)=\left\{\begin{array}{ll}e^{-\frac{1}{x^{2}}}, & x \neq 0 ; \\ 0, & x=0 .\end{array}\right.\) 则 \(f^{(2025)}(0)=\)
  \item[（8）] \(\iint_{x^{2}+y^{2} \leq 4} e^{-x^{2}-y^{2}} \mathrm{~d} x \mathrm{~d} y=\)
  \item[（9）] \(\oint_{L}(y-z) \mathrm{d} x+(z-x) \mathrm{d} y+(x-y) \mathrm{d} z=\) \(\_\_\_\_\)，其中 \(L\) 为柱面 \(x^{2}+y^{2}=R^{2}\) 与平面 \(\frac{x}{R}+\frac{y}{H}=\) \(1(R, H>0)\) 的交线，从 \(x\) 轴看沿逆时针方向．
  \item[（10）] \(\int_{0}^{+\infty} \frac{\sin (\sqrt{2} x)}{x} \mathrm{~d} x=\)
\end{itemize}

\subsection{题目2}
2．（15 分）证明 \(\left\{x_{n}\right\}\) 收敛，其中 \(x_{n}=\sum_{k=1}^{n} \frac{1}{k}-\ln n\) ．

\subsection{题目3}
3．（15 分）设 \(f(x)\) 在 \([0,+\infty)\) 上可导，且 \(f^{\prime}(x)\) 在 \([0,+\infty)\) 上连续，若 \(\lim _{x \rightarrow+\infty}\left[2 f(x)+f^{\prime}(x)\right]=0\) ，证明：

\[
\lim _{x \rightarrow+\infty} f(x)=0
\]

\subsection{题目4}
4．（15 分）设 \(f(x)\) 在 \([0,1]\) 上连续，在 \((0,1)\) 上可导，且 \(f(1)=2 \int_{0}^{\frac{1}{2}} x^{2} f(x) \mathrm{d} x\) ．证明：存在 \(\eta \in(0,1)\) ，使得 \(2 f(\eta)+\eta f^{\prime}(\eta)=0\) ．

\subsection{题目5}
5．（15 分）设 \(f(x)\) 在 \([0,2 \pi]\) 上无穷次可导，且 \(f(0)=f(2 \pi), \int_{0}^{2 \pi} f(x) \mathrm{d} x=0\) ．证明：

\[
\int_{0}^{2 \pi} f^{2}(x) \mathrm{d} x \leq \int_{0}^{2 \pi}\left[f^{\prime}(x)\right]^{2} \mathrm{~d} x
\]

\subsection{题目6}
6．（20 分）设 \(f(x)\) 在［ 0,1 ）上连续，证明：\(f(x)\) 在［0，1）上一致连续的充要条件是存在 \([0,1]\) 上的连续函数 \(g(x)\) ，当 \(x \in[0,1)\) 时，有 \(g(x)=f(x)\) ．

\subsection{题目7}
7．（20 分）解答如下问题：\\
（1）（10 分）设 \(f(x, y)\) 在 \(\mathbb{R}^{2}\) 上连续可微，且 \(\left|\frac{\partial f(x, y)}{\partial x}\right| \leq 1,\left|\frac{\partial f(x, y)}{\partial y}\right| \leq 1\) 对任意的 \((x, y) \in \mathbb{R}^{2}\) 成立．证明：对任意的 \(\left(x_{1}, y_{1}\right),\left(x_{2}, y_{2}\right) \in \mathbb{R}^{2}\) ，有

\[
\left|f\left(x_{1}, y_{1}\right)-f\left(x_{2}, y_{2}\right)\right|^{2} \leq 2\left(x_{1}-x_{2}\right)^{2}+2\left(y_{1}-y_{2}\right)^{2} .
\]

（2）（10 分）已知函数列 \(\left\{f_{n}(x, y)\right\}\) 在 \(\mathbb{R}^{2}\) 上连续可微，对任意的 \((x, y) \in \mathbb{R}^{2}\) 及正整数 \(n\) ，有

\[
\left|\frac{\partial f_{n}(x, y)}{\partial x}\right| \leq 1,\left|\frac{\partial f_{n}(x, y)}{\partial y}\right| \leq 1
\]

且 \(\lim _{n \rightarrow \infty} f_{n}(x, y)=f(x, y)\) 。证明：在 \(\mathbb{R}^{2}\) 的任意有界闭集中，\(\left\{f_{n}(x, y)\right\}\) 一致收敛于 \(f(x, y)\) 。
\end{document}
